% Volume VIII: Computational Neuroscience Mathematics
% Continuity note for Chapter 51.
% Previous dependencies: Chapters 13--18 provide probability, likelihood, uncertainty, and information; Chapters 31--34 provide statistical learning, GLMs, kernels, Gaussian processes, and latent variables; Chapters 46--50 provide spike-train models, neural coding, population geometry, circuits, and Bayesian decision models.
% New durable concepts: neural data modalities, observation units, preprocessing as a mathematical map, leakage, encoding and decoding models, likelihood-based fitting, train/test splits, cross-validation, AIC, BIC, held-out predictive comparison, posterior predictive checks, identifiability, uncertainty, construct validity, and ethical limits of neural modeling.
% Forward hooks: Chapter 52 compares biological and artificial systems using these analysis pipelines; Chapter 53 uses model comparison and data analysis inside integrated case studies; Chapter 54 summarizes the statistical and computational recipes.
% Notation commitments: Observations are indexed by m=1,\ldots,M; stimuli or covariates are x_m; neural responses are y_m or r_m; model parameters are theta; design matrices are X; response matrices are Y; training and test sets are I_train and I_test; log-likelihood is ell(theta;D).
\section{Chapter 51: Neural Data Analysis and Model Comparison}
\label{ch:neural-data-analysis-model-comparison}

Modern neuroscience produces many kinds of data: spike times, calcium fluorescence, local field potentials, EEG and MEG signals, fMRI time series, behavior videos, perturbation records, stimulus metadata, and trained model activations. These measurements do not become evidence merely because they are large. They become evidence when the observation unit is clear, preprocessing is controlled, models are fit on one part of the data and evaluated on another, and uncertainty is reported honestly.

\paragraph{Chapter problem.}
How can neural measurements be turned into analyzable mathematical objects, and how can competing models be compared without confusing prediction, explanation, identifiability, and ethical interpretation?

\paragraph{Section map.}
Section 51.1 classifies experimental data types and treats preprocessing as a mathematical transformation with possible leakage. Section 51.2 defines encoding and decoding models, likelihoods, train/test splits, and evaluation metrics. Section 51.3 develops predictive model comparison, AIC, BIC, posterior predictive checks, and identifiability. Section 51.4 discusses the epistemology and ethics of neural modeling: what claims the mathematics supports, what it does not support, and how uncertainty and privacy constraints enter.

\subsection{51.1 Experimental data types}

\paragraph{Guiding question.}
What exactly is one observation in a neural experiment, and how do modality, sampling, noise, and preprocessing determine the mathematical data object?

\paragraph{Beginner-facing intuition.}
A neural data set is not just an array. It is an array plus a measurement story. A spike count in a \(20\) ms bin, a calcium trace sampled at \(30\) Hz, an fMRI voxel time series sampled every two seconds, and a video-derived paw coordinate all have different time scales, noise models, and preprocessing needs. A model that ignores those differences may fit numbers while answering the wrong scientific question.

The first analysis decision is to define the observation unit. It may be one trial, one time bin, one neuron-time pair, one stimulus presentation, one behavioral episode, or one imaging frame. The second decision is to define the variables attached to that unit: neural response, stimulus, behavior, session, subject, experimental condition, and perturbation.

\paragraph{Formal setup.}
Let \(m=1,\ldots,M\) index observation units. A common supervised data object is
\[
D=\{(x_m,y_m,c_m):m=1,\ldots,M\}.
\]
Here \(x_m\in\mathcal X\) is a stimulus, task, or model covariate; \(y_m\in\mathcal Y\) is a neural or behavioral response; and \(c_m\in\mathcal C\) is context such as animal, session, trial type, or recording quality. For population recordings, one often stacks observations into matrices
\[
X\in\mathbb R^{M\times d},\qquad Y\in\mathbb R^{M\times N},
\]
where rows are observation units, \(d\) is the number of covariates, and \(N\) is the number of neurons, voxels, electrodes, or features.

A preprocessing pipeline is a map
\[
T:\mathcal D_{\mathrm{raw}}\to\mathcal D_{\mathrm{analyzed}}.
\]
Examples include filtering, motion correction, spike sorting, deconvolution, alignment to events, binning, normalization, artifact rejection, and dimensionality reduction. A learned preprocessing step has parameters \(\alpha\), so the analyzed data are \(T_\alpha(D_{\mathrm{raw}})\). To avoid leakage, \(\alpha\) must be estimated using only training data whenever the output will be evaluated on held-out data.

\begin{center}
\small
\renewcommand{\arraystretch}{1.18}
\begin{tabular}{p{0.17\linewidth}p{0.21\linewidth}p{0.23\linewidth}p{0.29\linewidth}}
\toprule
Modality & Typical response & Time scale & Analysis caveat \\
\midrule
Extracellular spikes & event times or binned counts & milliseconds & spike sorting, refractory violations, nonstationary firing rates \\
Calcium imaging & fluorescence trace or inferred events & tens to hundreds of milliseconds & slow indicator dynamics, motion correction, neuropil contamination \\
EEG/MEG/LFP & continuous field signal & milliseconds & filtering choices, volume conduction, spectral leakage \\
fMRI & voxel BOLD time series & seconds & hemodynamic response, autocorrelation, multiple comparisons \\
Behavior video & pose, speed, actions & frame rate dependent & tracking errors, label noise, coupling to neural timing \\
\bottomrule
\end{tabular}
\end{center}

\paragraph{Core mechanism: preprocessing and leakage.}
Suppose a scalar neural feature \(z_m\) is z-scored before fitting a decoder. If the mean and standard deviation are computed on all data,
\[
\tilde z_m=\frac{z_m-\bar z_{\mathrm{all}}}{s_{\mathrm{all}}},
\]
then test-set information enters the training representation. The safer split-aware version is
\[
\tilde z_m=
\frac{z_m-\bar z_{\mathrm{train}}}{s_{\mathrm{train}}},
\qquad m\in I_{\mathrm{train}}\cup I_{\mathrm{test}},
\]
where \(\bar z_{\mathrm{train}}\) and \(s_{\mathrm{train}}\) are estimated only from \(I_{\mathrm{train}}\). The test feature is transformed with training-set statistics. This is not pedantry: if a session has slow drift, using all trials to normalize can make future test trials partly visible during training.

\paragraph{Concrete example.}
Three trials have stimulus strengths \(x=(0,1,1)\) and two-neuron spike counts
\[
Y=
\begin{bmatrix}
1&0\\
3&1\\
5&2
\end{bmatrix}.
\]
If each row is one trial, then \(M=3\) and \(N=2\). If the first two trials are training data, the training mean for neuron 1 is \((1+3)/2=2\), not the all-data mean \(3\). A decoder evaluated on trial 3 should use the training mean \(2\) for normalization, because trial 3 is supposed to be held out.

\paragraph{Computational neuroscience example.}
In a reaching experiment, raw voltage traces may be spike-sorted into spike times, binned relative to movement onset, and converted to firing rates. Each preprocessing choice changes the response variable \(Y\). A \(5\) ms bin preserves timing but is sparse; a \(100\) ms bin is less sparse but hides fast dynamics. The right bin size depends on the scientific question and the model's assumptions.

\paragraph{AI or machine-learning example.}
The same leakage issue appears when standardizing image features or language-model activations before fitting a probe. If normalization, PCA, or feature selection is fit on the full data set before splitting, the held-out score is optimistic. The pipeline should be treated as part of the model.

\paragraph{Common misconceptions and failure modes.}
Raw data are not automatically more objective than preprocessed data; raw recordings contain artifacts and sampling constraints. Preprocessing is also not harmless, because it can remove signal or introduce leakage. Trial averaging improves signal-to-noise for condition means but can erase trial-to-trial structure. Finally, a modality's numerical resolution is not its scientific resolution: fMRI can have many voxels while still being slow relative to spiking dynamics.

\paragraph{Exercises.}
\begin{enumerate}[label=\textbf{51.1.\arabic*.}, leftmargin=*]
\item \textbf{Mechanical.} A recording has \(40\) trials, \(80\) time bins, and \(120\) neurons. If each trial-time pair is one observation and each neuron is a response coordinate, state \(M\), \(N\), and the shape of \(Y\).
\item \textbf{Proof or derivation.} Show that centering test data by the all-data mean makes the transformed test feature depend on the test feature itself. Write \(\bar z_{\mathrm{all}}\) in terms of training and test sums.
\item \textbf{Numerical/computational.} For training values \(2,4,6\) and test value \(10\), compute the training mean and the all-data mean. How would the z-scored test value differ if both standard deviations were \(2\)?
\item \textbf{Interpretation.} Why might a \(10\) ms spike-count bin and a \(500\) ms fMRI sample be inappropriate inputs to the same time-resolved model without additional assumptions?
\item \textbf{Misconception/debugging.} A student says, ``Preprocessing happens before modeling, so it cannot affect model comparison.'' Explain why learned preprocessing parameters must be included in the modeling pipeline.
\end{enumerate}

\subsection{51.2 Encoding and decoding models}

\paragraph{Guiding question.}
What is the difference between predicting neural responses from stimuli and predicting stimuli or behavior from neural responses?

\paragraph{Beginner-facing intuition.}
Encoding and decoding reverse the direction of prediction. An encoding model asks: given a stimulus, task variable, or artificial-network feature, what neural response should we expect? A decoding model asks: given neural activity, what stimulus, action, or latent variable can we estimate? Both can be useful, but they answer different scientific questions.

An accurate decoder does not prove that the brain uses that decoder. It proves that information about the target variable is present in the measured response under the chosen data distribution and metric. An accurate encoding model does not prove mechanism either, but it can show that a proposed feature set predicts neural responses.

\paragraph{Formal setup.}
An encoding model specifies a conditional distribution
\[
p_\theta(y\mid x,c),
\]
where \(x\) is a stimulus or covariate, \(c\) is context, \(y\) is neural response, and \(\theta\in\Theta\) is a parameter. The log-likelihood on a data set \(D\) is
\[
\ell(\theta;D)=\sum_{m=1}^M \log p_\theta(y_m\mid x_m,c_m).
\]
A decoding model specifies either a conditional distribution
\[
q_\phi(x\mid y,c)
\]
or a deterministic estimator
\[
g_\phi:\mathcal Y\times\mathcal C\to\mathcal X.
\]
For regression targets, a common decoder minimizes squared error
\[
\sum_{m\in I_{\mathrm{train}}}\|x_m-g_\phi(y_m,c_m)\|_2^2.
\]
For classification targets, a common decoder minimizes cross-entropy.

Train/test discipline separates fitting from evaluation:
\[
I_{\mathrm{train}}\cap I_{\mathrm{test}}=\emptyset.
\]
Parameters are chosen using \(I_{\mathrm{train}}\), and performance is reported on \(I_{\mathrm{test}}\).

\paragraph{Core mechanism: Gaussian encoding regression.}
Let \(x_m\in\mathbb R^d\) and scalar response \(y_m\in\mathbb R\). A linear-Gaussian encoding model is
\[
y_m=x_m^\top\beta+\epsilon_m,\qquad \epsilon_m\sim\mathcal N(0,\sigma^2).
\]
The log-likelihood is
\[
\ell(\beta,\sigma^2)
=-\frac{M}{2}\log(2\pi\sigma^2)
-\frac{1}{2\sigma^2}\sum_{m=1}^M (y_m-x_m^\top\beta)^2.
\]
For fixed \(\sigma^2\), maximizing likelihood is the same as minimizing squared error. If \(X\in\mathbb R^{M\times d}\) has rows \(x_m^\top\) and \(y\in\mathbb R^M\), the least-squares normal equations are
\[
X^\top X\widehat\beta=X^\top y.
\]
When \(X^\top X\) is invertible,
\[
\widehat\beta=(X^\top X)^{-1}X^\top y.
\]
Thus a familiar regression fit is also a likelihood fit under an explicit noise model.

\paragraph{Concrete example.}
Suppose one stimulus feature \(x\) and one neuron's response \(y\) are observed on three training trials:
\[
(x,y)=(0,1),(1,3),(2,5).
\]
The linear model \(y=\beta_0+\beta_1x\) fits exactly with \(\widehat\beta_0=1\) and \(\widehat\beta_1=2\). On a held-out trial with \(x=3\), the encoding prediction is \(7\). If the observed response is \(6\), the squared prediction error is \(1\).

\paragraph{Computational neuroscience example.}
For a visual neuron, an encoding model may use image features \(x_m\) to predict spike counts \(y_m\). A Poisson GLM might set
\[
\lambda_m=\exp(\beta_0+x_m^\top\beta),\qquad
Y_m\mid x_m\sim\operatorname{Poisson}(\lambda_m).
\]
For decoding, the response vector from many neurons may be used to predict reach direction or stimulus orientation. Encoding tells how stimulus features drive measured neurons; decoding tells what can be inferred from the population response.

\paragraph{AI or machine-learning example.}
In systems neuroscience, activations from a trained CNN or transformer can serve as \(x_m\) in an encoding model for brain responses. Conversely, a neural decoder can be a classifier trained on response vectors. The same train/test and likelihood rules apply to artificial activations and biological recordings.

\paragraph{Common misconceptions and failure modes.}
Encoding and decoding are not inverses unless a full probabilistic model, prior, and loss are specified. High decoding accuracy can be caused by confounds, such as movement or arousal, rather than the intended stimulus variable. High training accuracy is not evidence of generalization. A linear model can be scientifically useful even when the brain is nonlinear, because it tests a restricted predictive hypothesis.

\paragraph{Exercises.}
\begin{enumerate}[label=\textbf{51.2.\arabic*.}, leftmargin=*]
\item \textbf{Mechanical.} State whether each model is encoding or decoding: \(p_\theta(r\mid s)\), \(q_\phi(s\mid r)\), \(g(r)=\widehat{\text{reach direction}}\), and \(\lambda(s)=\exp(\beta_0+\beta_1s)\).
\item \textbf{Proof or derivation.} Derive why maximizing the Gaussian linear-model likelihood for fixed \(\sigma^2\) is equivalent to minimizing the sum of squared residuals.
\item \textbf{Numerical/computational.} For \(y=1+2x\), compute predictions at \(x=0,1,2,3\). If the held-out responses are \(1,4,4,8\), compute the mean squared error.
\item \textbf{Interpretation.} What scientific question is answered by an encoding model that predicts V1 responses from image features? What different question is answered by a decoder that predicts image category from V1 activity?
\item \textbf{Misconception/debugging.} A decoder predicts stimulus identity well from neural activity. Explain why this does not prove that downstream brain areas use the same decoder.
\end{enumerate}

\subsection{51.3 Model comparison and identifiability}

\paragraph{Guiding question.}
When two models fit the same data, how do we decide which one predicts better, which one is simpler, and whether its parameters are actually identifiable?

\paragraph{Beginner-facing intuition.}
Model comparison is not the same as picking the model with the smallest training error. Flexible models can fit noise. Simple models can miss real structure. The safest default is predictive comparison on data not used for fitting. Likelihood-based criteria such as AIC and BIC add penalties for parameter count. Posterior predictive checks ask whether simulated data from the fitted model look like real data in targeted ways.

Identifiability asks a different question. A parameter is identifiable if different parameter values imply different observable distributions. If two parameter settings make the same predictions for every possible observation, no amount of data from that experiment can distinguish them.

\paragraph{Formal setup.}
Let \(M_j\) be a model class with parameter \(\theta_j\in\Theta_j\) and likelihood \(p_j(D\mid\theta_j)\). The maximum log-likelihood is
\[
\ell_j(\widehat\theta_j;D)=\max_{\theta_j\in\Theta_j}\log p_j(D\mid\theta_j).
\]
If \(k_j\) is the number of fitted parameters and \(n\) is the number of independent observation units, two common criteria are
\[
\boxed{\operatorname{AIC}_j=-2\ell_j(\widehat\theta_j;D)+2k_j}
\]
and
\[
\boxed{\operatorname{BIC}_j=-2\ell_j(\widehat\theta_j;D)+k_j\log n.}
\]
Lower values are preferred by these criteria.

Held-out predictive comparison uses
\[
\operatorname{LL}_{\mathrm{test}}(M_j)
=\sum_{m\in I_{\mathrm{test}}}\log p_j(y_m\mid x_m,\widehat\theta_j).
\]
Cross-validation averages this quantity over multiple train/test splits.

Identifiability is a property of the map
\[
\theta\mapsto P_\theta.
\]
A parameter \(\theta\) is globally identifiable if
\[
P_\theta=P_{\theta'}\quad\Longrightarrow\quad \theta=\theta'.
\]

\paragraph{Core mechanism: predictive comparison and posterior predictive checks.}
Held-out likelihood evaluates the probability assigned to data the model did not see during fitting. If two models have similar training likelihood but one gives much higher held-out likelihood, the latter generalizes better under the chosen split.

A posterior predictive check, in Bayesian form, draws parameters and replicated data:
\[
\theta^{(s)}\sim p(\theta\mid D),\qquad
D_{\mathrm{rep}}^{(s)}\sim p(D_{\mathrm{rep}}\mid \theta^{(s)}).
\]
Then a statistic \(T(D)\), such as spike-count variance, autocorrelation, peak response time, or decoding accuracy, is compared with \(T(D_{\mathrm{rep}}^{(s)})\). The goal is not to prove the model true. The goal is to find important ways in which it fails.

\paragraph{Concrete example.}
Model A has maximum log-likelihood \(-120\) and \(k=3\) parameters. Model B has maximum log-likelihood \(-112\) and \(k=10\). Their AIC values are
\[
\operatorname{AIC}_A=240+6=246,\qquad
\operatorname{AIC}_B=224+20=244.
\]
AIC slightly prefers B. If \(n=100\), their BIC values are
\[
\operatorname{BIC}_A=240+3\log 100\approx253.8,
\]
\[
\operatorname{BIC}_B=224+10\log 100\approx270.1.
\]
BIC prefers A because it penalizes the extra parameters more strongly.

For non-identifiability, consider the model
\[
y=(ab)x.
\]
Only the product \(ab\) affects predictions. The parameter pairs \((a,b)=(2,3)\) and \((1,6)\) produce the same output for every \(x\). Therefore \(a\) and \(b\) are not separately identifiable from this experiment.

\paragraph{Computational neuroscience example.}
Two spike-train models may have similar training likelihood: a Poisson model with stimulus drive only and a GLM with both stimulus and spike-history filters. Held-out likelihood can test whether the history term predicts new spikes. Posterior predictive checks can test whether simulated spike trains reproduce refractory periods, burstiness, and condition-dependent variability.

\paragraph{AI or machine-learning example.}
Deep networks are often compared by validation loss, calibration, robustness, and out-of-distribution performance rather than training loss alone. Identifiability issues also arise: permuting hidden units in a neural network can produce the same function, so individual weights are not always interpretable even when predictions are stable.

\paragraph{Common misconceptions and failure modes.}
AIC and BIC are not universal truth meters; they rely on assumptions and approximations. Cross-validation estimates predictive performance for a particular data-splitting scheme, so leakage or correlated trials can make it optimistic. Identifiability is not the same as low uncertainty: a parameter can be identifiable but poorly estimated with small data, or non-identifiable even with infinite data. A posterior predictive check can reveal model failure without specifying the best alternative.

\paragraph{Exercises.}
\begin{enumerate}[label=\textbf{51.3.\arabic*.}, leftmargin=*]
\item \textbf{Mechanical.} Compute AIC for a model with log-likelihood \(-80\) and \(k=5\). Compute BIC when \(n=50\).
\item \textbf{Proof or derivation.} Show that in the model \(y=(ab)x\), the parameter pair \((a,b)\) is not globally identifiable by constructing infinitely many pairs with the same product.
\item \textbf{Numerical/computational.} Two models have held-out log-likelihoods \(-31.2\) and \(-29.5\) on the same test trials. Which assigns higher predictive probability to the test data? What is the log-likelihood difference?
\item \textbf{Interpretation.} A fitted Poisson model matches mean spike counts but simulated data have much smaller variance than real data. What posterior predictive failure does this suggest?
\item \textbf{Misconception/debugging.} A student says, ``The model with the most parameters is best because it can fit the data most closely.'' Explain why held-out prediction and parameter penalties matter.
\end{enumerate}

\subsection{51.4 Ethics and epistemology of neural modeling}

\paragraph{Guiding question.}
What can a neural model justify us in claiming, and what responsibilities arise when neural data are private, noisy, biased, or mechanistically ambiguous?

\paragraph{Beginner-facing intuition.}
Neural models sit between measurement and claim. A classifier may predict disease status from brain data. A decoder may infer a stimulus from spikes. A large artificial network may predict responses in cortex. These results can be useful, but the mathematical output is not automatically a mind reading, a causal explanation, or a clinical decision.

Epistemology asks what kind of knowledge the analysis supports. Ethics asks whether the data collection, modeling, reporting, and use respect participants, animals, patients, and affected communities. Good mathematics helps both because it makes assumptions, uncertainty, and failure modes explicit.

\paragraph{Formal setup.}
Let \(D\) be a data set sampled from an unknown process \(P\). A scientific claim often concerns an estimand
\[
\psi(P),
\]
such as a tuning strength, decoding accuracy on a target population, causal effect, or model-predicted response. An estimator is a function of data,
\[
\widehat\psi=\widehat\psi(D).
\]
Uncertainty is represented by an interval, posterior distribution, bootstrap distribution, standard error, or predictive distribution. A claim should state the target population, observation unit, model class, and uncertainty procedure.

Construct validity asks whether \(\psi(P)\) matches the intended concept. For example, high decoding accuracy for a stimulus label is an estimand about predictability from measured responses; it is not automatically an estimand of conscious perception or causal neural representation.

Privacy risk can be formalized by the information a released statistic \(A(D)\) carries about an individual record \(D_i\). Differential privacy, one rigorous framework, requires roughly that changing one individual's data changes the distribution of the released output only a little:
\[
P(A(D)\in S)\leq e^\varepsilon P(A(D')\in S)
\]
for neighboring data sets \(D,D'\) differing in one record and for output sets \(S\). This course does not develop privacy theory, but the inequality shows that privacy is a mathematical constraint, not only a policy slogan.

\paragraph{Core mechanism: separating prediction, mechanism, and causation.}
A predictive model estimates conditional structure such as
\[
p(y\mid x).
\]
A mechanistic model proposes internal variables and equations that could generate the data. A causal model makes claims about interventions, such as
\[
P(Y\mid \operatorname{do}(X=x)).
\]
These are different objects. A decoder with high \(P(\widehat X=X)\) estimates predictive information in the measurement. It does not by itself identify an intervention effect or a biological algorithm.

Uncertainty should travel with the claim. If a held-out decoding accuracy is \(\widehat a=0.82\) on \(n=100\) independent trials, an approximate binomial standard error is
\[
\sqrt{\frac{\widehat a(1-\widehat a)}{n}}
=\sqrt{\frac{0.82(0.18)}{100}}
\approx0.038.
\]
The point estimate is incomplete without the sampling context.

\paragraph{Concrete example.}
Suppose a model predicts a binary disease label from fMRI with \(82\) correct predictions out of \(100\) held-out subjects. The accuracy is \(0.82\). This supports a predictive claim for subjects drawn like the test set. It does not prove that the model has found the disease mechanism, that it generalizes to different scanners or populations, or that it should be used clinically without calibration, external validation, and harm analysis.

\paragraph{Computational neuroscience example.}
A perturbation experiment may show that stimulating a population changes behavior. This is stronger causal evidence than an observational correlation, but interpretation still depends on stimulation specificity, off-target effects, behavioral confounds, and model assumptions. Ethical analysis also includes animal welfare and the necessity of the perturbation.

\paragraph{AI or machine-learning example.}
Brain-computer interfaces and neural decoders can improve communication or motor control, but they also raise privacy and autonomy concerns. A model that decodes intended movement under consented laboratory conditions should not be described as a general-purpose thought decoder. The training distribution, consent boundary, uncertainty, and failure modes are part of the scientific claim.

\paragraph{Common misconceptions and failure modes.}
Predictive success is not the same as explanation. Statistical significance is not practical or ethical significance. A confidence interval does not account for every source of bias, such as sampling bias, measurement artifacts, or construct mismatch. Neural data are not inherently anonymous; high-dimensional biological and behavioral measurements can identify people or reveal sensitive attributes. Finally, cautious wording is not weakness; it is part of making the claim mathematically and ethically valid.

\paragraph{Exercises.}
\begin{enumerate}[label=\textbf{51.4.\arabic*.}, leftmargin=*]
\item \textbf{Mechanical.} Classify each claim as predictive, mechanistic, or causal: ``the decoder predicts reach direction,'' ``the circuit equation produces oscillations,'' and ``silencing this population reduces choices to the right.''
\item \textbf{Proof or derivation.} Derive the binomial standard error formula \(\sqrt{\widehat a(1-\widehat a)/n}\) by using the variance of a Bernoulli random variable and the variance of a sample mean.
\item \textbf{Numerical/computational.} A decoder gets \(45\) out of \(60\) test trials correct. Compute accuracy and the approximate binomial standard error.
\item \textbf{Interpretation.} Why does high decoding accuracy from neural data not automatically imply that the decoded variable is consciously represented?
\item \textbf{Misconception/debugging.} A report says, ``The model reads private thoughts because it predicts one task variable above chance.'' Rewrite the claim in a more mathematically defensible form.
\end{enumerate}

\paragraph{Closing bridge.}
This chapter treated neural data analysis as disciplined inference: define the observation, control preprocessing, fit models on training data, compare them on held-out or likelihood-based criteria, check failures, and state claims with uncertainty. Chapter 52 uses this discipline to compare biological and artificial intelligence without shallow analogy. Chapter 53 then integrates the course's main mathematical tools through case studies, and Chapter 54 collects the reference formulas needed to carry out those analyses.

\clearpage
